\chapter{撤离与着陆}
\chapterkana{撤退と着陸}
\epigraphbox{“颛顼受之，乃命南正重司天以属神，命火正黎司地以属民，使复旧常，无相侵渎，是谓绝地天通。”}{——《国语·楚语下》}

\kaishou{在城里撑了一千七百余年之后，地下的热对流先一步凉了下来。\par{}十二万人分六批回到地面，留下一座四千七百一十三米长的空壳撞进黄土。\par{}最后一道垂直通道在四十米处截断，断面旁只刻了两个中文：「勿挖」。}

\section{六批撤离}

单靠净秽山内越攒越厚的活性污秽，还不足以令一座支撑了近两千年的城市在某一年突然全体撤空；真正按下撤离开关的，是能源链路的衰竭。积累是远因，链路衰微才是压垮系统的近因。

这座城市的空气、温湿、定向风以及压制净秽山的能耗，从正面图纸上看不到本地发电单元。第一章引过的俄科院第 117 全宗第 2187 号卷宗第 19 叶（1:200 纵剖面）上，静海七号中央枢纽底层画着一组向下延伸的「热流耦合受能阵」，通篇没有自备反应堆，第 20 叶只注了「第 4 约束段，检修周期 4000 h」。这组无源受能结构究竟接向哪里，要对照另一份冷战时期的地月遥测简报（档案 \texttt{ФВ-1968-29}）才接得上头。

\begin{archive}[内部参考]{档案 № ФВ-1968-29}
一九六八年深地热流与轨道耦合观测备忘（苏联科学院远东地球物理站移交）：\\
1967—1968 年间在欧亚大陆三口深钻井测得周期约为 4000 小时的微弱热流脉动，其相位锁定角与静海西缘地下空腔残余谐振频段吻合。对比 1961 年月面工程图纸第 19 叶耦合阵倾角，该链路在结构上指向以地球地核热对流作远端供能端、由月面耦合阵接收的定向能传输假说。
\end{archive}

在把这座城看作靠地球行星级热机远程受能的工程推论下，能源衰退的脉络便清晰了许多：据俄科院第 117 全宗档案第 2187 号卷宗所存数据，并与《广寒宫残简》简二十七对校，大撤离末期约二百年间，地核热对流效率先慢后快地累计下降了千分之三十一（即 3.1\%）。\footnote{第 117 全宗第 2187 号卷宗参证《广寒宫残简》简二十七「三十一分」之数。此处「分」为千分率，两处出处互不相隶，数字完全合拍。}起初只是外环气压与水路偶尔吃紧，到约公元前四五五年前后，山内积压了一千五百余年的活性污秽逼近临界，链路冷却迹象也压抑不住，月之都便开始分批向地面遣送居民。

核对六批人数与落点的主证物，是编入基金会目录的《撤离签符与登船核对册》（档案 \texttt{ФВ-1970-18}），内收四片带编号的青铜签符与一份残缺的登船总录。签符长约七厘米，每片正面阴刻人名，背面空着，与《广寒宫残简》所载「不交回、不补发」的登船凭证形制相合。

\begin{archive}[内部参考]{档案 № ФВ-1970-18}
登船总录残页两行：「前五次合五万五千，依岁序遣。」「末次未入原籍册者千，实载六万五千，与原册十一万九千差千。」
\end{archive}

把 \texttt{ФВ-1970-18} 与前文用到的《章程总目》题头、名籍末枚，以及《广寒宫残简》简十九至三十的清点记录并列一处，整场撤离的人口账目便显出了全貌，连那些至今对不齐的尾数也一目了然（见表 \ref{tab:evacuation}）。其中有实载之数，也有在册之数，彼此不能互相代算。

\begin{table}[htbp]
  \centering
  \caption{六批撤离与全城人口口径对照}
  \label{tab:evacuation}
  \footnotesize
  \begin{tabular}{@{}p{26mm}p{30mm}p{18mm}p{48mm}@{}}
    \toprule
    批次／口径 & 年代 & 数目（人） & 备注 \\
    \midrule
    开工计口 & 船始建时 & 312{,}000 & 静海七号与背面城设计额定容量 \\
    完工登船 & 建城毕时 & 150{,}000 & 题头「完工之日，十五万」 \\
    兔乱折损 & 约前 1980 年 & 22{,}000 & 月之民死亡数（另余 8{,}000 无着落） \\
    第一至五批 & 约前 455—前 255 年 & 55{,}000 & 计划内每约四十五至五十年送出一批（6k、8k、10k、13k、18k） \\
    第六批（实载） & 前 221 年 & 65{,}000 & 生态崩溃后紧急撤离；六批实载总计 120{,}000 \\
    撤离在册 & 前 221 年 & 119{,}000 & 名籍末枚所载总数，较实载少 1{,}000 \\
    \bottomrule
  \end{tabular}
  \tabsource{合参《章程总目》题头、名籍卷、\texttt{ФВ-1970-18} 与《月之都——〈广寒宫残简〉辑佚与考释》简十九至三十。}
\end{table}

前五批约每四十五至五十年按岁计计划发出一批（依序为六千、八千、一万、一万三千、一万八千人），在城内维持与垂直运输体系尚可运转时分乘常规登陆舱段降落，散入中原各地；其后仅过三十余年，到公元前二二一年，地热加速衰减令定向能链路跌穿临界阈值，城市气路与水路彻底断绝，系统全面崩溃，剩下六万五千人已无从按既有计划从容转运，只得依托重型着陆载具仓促起运。两笔差额依然没有强行抹平的余地：一笔是完工时的十五万扣除兔乱二万二千及六批实载十二万后尚余的八千人；另一笔是实载十二万与在册十一万九千之间多出的一千人。至于前二二一年撤离时唯一未随行者，便是第九章一直追踪的那位隔离舱服药者：她体内不断产生污秽，若随船降落或共处一舱，其余人立时就会沾染变化。撤离人员切断隔离设施供电后封门离去，此后再没有任何档案记下她的生死或去向。自前二二一年那一日起，整座月球上再无其余月之民。

\section{「昆仑山号」}

撤离所用的工程载具留下了两件互相独立的地面合金证物。第一件是远东委员会于一九五九至一九六二年在蒙古肯特省、内蒙古鄂尔多斯高原及赤塔州一带采集的样品分析卷宗，编号 \texttt{ФВ-1963-07}，其中最具决定性的是鄂尔多斯样品 \textsc{э}-14。

\begin{archive}[绝密／СОВЕРШЕННО СЕКРЕТНО]{档案 № ФВ-1963-07}
本卷宗系一九六三年九月由苏联科学院远东委员会第四组移交，共 217 叶。第 88 叶记载，鄂尔多斯样品 \textsc{э}-14 为一枚重 41.6 克的片状物，经分析为 \textsc{nb}—\textsc{ta}—\textsc{w} 三元合金。其铌同位素比值 \textsc{nb}-94/\textsc{nb}-93 为 0.0172，与地球物质标准值 0.0170 在误差范围之内一致；但 \textsc{ta}-180 的丰度较地球标准值高出 2.3 倍，这一差异无法用区域地质背景加以解释。\\
第 89 叶页边存有巴兰尼科夫本人铅笔批注：「坦塔洛斯之泪。他们把它带下来了。」其下另有一行被墨水涂黑，红外成像可辨出「重量不对」四字。
\end{archive}

第二件碎块出土于豫鲁交界处相当于秦汉初年淤积层，编号 \mbox{\texttt{ФВ-1974-12}：}全块长十八点四厘米、重二点三六公斤，表面覆着一层平均厚度约一点七毫米的黑褐色气动熔蚀壳，内部基体却呈现出含铌、钛的玄武质致密合金结构，微量元素配比与第六章所见的月面断钻头附着物高度相近。\footnote{标本采集于山东鄄城旧河床探方，同层伴出秦末半两钱范碎片两件。}按年代与地点比对，它恰好对应《史记·秦始皇本纪》三十六年（前二一一年）「有坠星下东郡，至地为石」的那条记录，极可能是高空载具或伴随碎片解体后迟落东侧的一枚残片。\footnote{《史记》卷六《秦始皇本纪》，中华书局 1959 年点校本，第 259 页。}

这两枚带着月面冶炼配比的地面残片，连同第 117 全宗的航迹图纸，勾出了公元前二二一年那套紧急撤离载具的分工：母船在档案中称为「昆仑山号」，本身留在地球轨道上（即第二章讨论过的月球主体，正面静海七号仅以约两成的实际占有率空置封存）；真正脱离轨道再入大气层的着陆舱组，则总称为「玉山」。\footnote{轨道与月壳参数的工程复原亦见巴兰尼科夫《月面设施复原图集》第 7 页，圣彼得堡，2016 年。}「玉山」由一座长达四千七百一十三米的无人主舱加上十二个乘员舱组成：十二个乘员舱在再入中途分离减速，分散降落于豫西与关中一带，无人主舱则一路坠入黄河流域。

\section{「玉山」再入}

关于「玉山」究竟如何穿过大气层，俄科院第 117 全宗第 2187 号卷宗第 41 叶保留了一条包含两条候选方案的速度衰减曲线（见图 \ref{fig:traj}）。两条曲线代表了两种截然不同的再入构型。

\begin{figure}[htbp]
\centering
\begin{tikzpicture}[scale=0.95, font=\small]
  \begin{axis}[
    width=12.6cm, height=5.8cm,
    xlabel={飞行时间（小时）}, ylabel={},
    xmin=0, xmax=264, ymin=0, ymax=12.5,
    axis lines=left, tick align=outside,
    xtick={0,48,96,144,192,240},
    ytick={0,2,4,6,8,10,12},
    every axis label/.style={font=\footnotesize},
    tick label style={font=\scriptsize},
    grid=major, grid style={rulegrey!45, line width=0.2pt},
    ]
    \addplot[line width=0.9pt, domain=0:264, samples=120]
      {12.1*exp(-x/90)+0.55};
    \addplot[line width=0.5pt, dashed, domain=0:264, samples=120]
      {9.4*exp(-x/74)+0.42};
    \node[font=\scriptsize, anchor=west] at (axis cs:150,6.4) {方案 A（静海直降）};
    \node[font=\scriptsize, anchor=west] at (axis cs:150,3.1) {方案 B（近地制动）};
    \draw[line width=0.3pt, dashed] (axis cs:0,0.55) -- (axis cs:264,0.55);
    \node[font=\scriptsize, anchor=north west] at (axis cs:14,0.42) {乘员舱末速 0.55 km/s};
    \node[rotate=90, font=\footnotesize] at (rel axis cs:-0.115,0.5) {速度（km/s）};
  \end{axis}
\end{tikzpicture}
\caption{「玉山」乘员舱再入速度衰减曲线（两条候选方案）}
\label{fig:traj}
\figsource{据第 117 全宗第 2187 号卷宗第 41 叶数据点重绘。原图纵轴无标注，数值由俄文标记「в.\,е.」换算而来，换算系数见 Баранников 2018 年书第 244 页。}
\end{figure}

方案 A 保持较高速度穿越大气，直接以四点七公里长的无人主舱本身充当气动减速面，一路烧穿落地而不作完整缓冲，十二个乘员舱则在中途脱离、各自减速；方案 B 在近地轨道先行完成制动，以较低速度温和进入，但要求所有着陆舱均具备完整的热防护结构。档案在两条曲线旁各留有一处极简批注：方案 A 旁写「已执行」，方案 B 旁写「未执行」。按当时地核供能骤断、系统已难维持近地制动的仓促处境来看，放弃方案 B 而改行以主舱硬扛减速的方案 A，反倒更符合现实工程约束。

既然执行的是方案 A，其力学后果便是一座长度超过四点七公里、质量逾亿吨的主体在公元前二二一年高超音速坠入黄河流域，具体参数复原列于表 \ref{tab:landing}。

\begin{table}[htbp]
  \centering
  \caption{「玉山」主舱降落参数复原（方案 A）}
  \label{tab:landing}
  \footnotesize
  \begin{tabular}{@{}llr@{}}
    \toprule
    参数 & 数值 & 不确定度 \\
    \midrule
    主舱长度 & 4{,}713 м & ± 12 м \\
    再入质量 & 1.1 × 10\textsuperscript{8} t & — \\
    触地速度 & 12.65 km/s & ± 0.2 km/s \\
    大气进入角 & 6.2° & ± 0.3° \\
    峰值热流 & 1.47 × 10\textsuperscript{4} kW/m² & ± 4\% \\
    落点 & 34.6°N, 110.4°E & ± 40 км \\
    撞击当量 & 2.1 × 10\textsuperscript{3} Mt & — \\
    抛射物总量 & 9.4 × 10\textsuperscript{12} t & — \\
    \bottomrule
  \end{tabular}
  \tabsource{据巴兰尼科夫《月面设施复原图集》第 251 页及第 117 全宗档案。落点坐标 34.6°N、110.4°E 位于今河南省灵宝市以西，处于仰韶文化西缘与二里头文化之间的过渡地带。}
\end{table}

表中的「抛射物总量」一栏，牵涉到中国北方第四纪晚期地层的一桩公案。一次当量为两千一百兆吨的撞击，会把约九万四千亿吨的物质抛入平流层；这些细颗粒物质在随后约三百年间缓慢沉降，落于黄土高原者平均约为每平方米十五吨，压实后折合十余米，正是今日所见部分次生黄土层的成因。常规地质学将黄土高原的沉积归为风力搬运，这一解释本身并不违背观测，只是在公元前二二一年之后的那数百年间，将如此巨量的粉尘抛送上空的原始动力不是季风，而是这次主舱撞击的冲击波，风力只是后续的再分配者。

\section{「绝地天通」}

十二乘员舱散落于豫西与关中之后，六万五千名末批撤离者与此前五批先期落地者合流，共计十二万人。这批长期生活在低重力、封闭气压与人工光照下的居民一踏进地面一点二个大气压、百分之二十一氧分压且尘埃弥漫的自然大气，除了承受剧烈的环境与生理负荷之外，在工程与居住方式上也立刻采取了保守的收缩措施：他们没有另起炉灶新建大城，而是就近接管并改造了早已存在数百上千年的仰韶、龙山与二里头旧址。

考古记录在多处早期遗址原生居住面以下发现了与原建筑年代不符的改造痕迹：在西安半坡，两座废弃房址的居住面被重新下挖约四十厘米，坑壁另加夯筑一层；在河南偃师二里头第三发掘区，宫殿基址年代远早于前二二一年约一千三百年，柱洞底部却各有一圈高温环状烧结面，其烧成温度远远高于当地炊煮或制陶遗迹。由于这些异常痕迹全部位于原生居住面之下，难以简单归入盗掘或表土扰动；结合当时落地人群避尘避风的实际需要来看，这些动作无非是加深地穴、封堵风口与将旧窑洞改作滤室，利用既有的地下结构勉强维持过渡时期的生存环境。

与旧址改造同样重要的，是公元前二二一年撤离收尾时对地月垂直运输通道的彻底破坏。在传世文献《国语·楚语下》里，观射父曾追述上古一句极有名的记载：「及少皞之衰也，九黎乱德……颛顼受之，乃命南正重司天以属神，命火正黎司地以属民，使复旧常，无相侵渎，是谓绝地天通。」历代注疏多将「绝地天通」解作王室垄断祭祀与通天权柄的宗教改革，但在文本之外，真正落在工程层面的实体线索却藏在莫斯科档案与二里头的柱洞底部。

第 117 全宗第 2187 号卷宗第 121 叶是一份用打字机誊抄的「设施退役清单」，共十九项，其中第七项写得极为直白：「垂直运输装置（电梯）已拆除，支柱于四十米高度截断，截断处标记『勿挖』。」且档案特别注明，该警告标记所用的文字正是中文。在地面与它恰好形成互校的，则是一九七一年在河南偃师二里头遗址发掘出的一组地下异常柱洞群（档案 \texttt{ФВ-1971-22}）。

\begin{archive}[绝密／СОВЕРШЕННО СЕКРЕТНО]{档案 № ФВ-1971-22}
一九七一年，中国科学院考古研究所在河南偃师二里头遗址第三发掘区，于 T-14 探方距地表 4.2 米处发现一处柱洞群，共九个柱洞，呈正圆形排列，直径 3.1 米。洞内填充物为经高温烧结的土块，烧结温度经热释光测定为 1{,}240 ℃ ± 60 ℃。\\
同层出土陶片七件，其中一件（编号 71YLT14:6）内壁有朱书符号两个，经红外成像可辨为「天」「梯」二字。字迹草率，笔画末端有拖曳痕迹，书写时似甚仓促。\footnote{见《河南偃师二里头遗址三区发掘简报》，《考古》1974 年第 5 期，第 268—279 页。简报第 273 页原载「柱洞底部发现不明金属残片，已送有关部门鉴定」一句，其后重印本中已予删去，至今未见公开鉴定结论。}
\end{archive}

倘若把「勿挖」的中文警示、四十米处截断的垂直支柱清单、九个呈正圆形排列的高温烧结柱洞以及那块仓促写就「天」「梯」二字的陶片摆在一张桌上，便能明白后世神话究竟把什么事情压缩成了「绝地天通」。在失去月面供能与大型冶炼供应链之后，地上仅存的少量落难者必须防止后来不知情的本土人群往下掘开残留的深埋通道或高压设施，而在技术语汇行将散尽之际，最经济有效的编码方式便是把工程禁令写入仪式与传说：把「勿挖」写成「不可亵渎」，把被截断的轨道垂直运输装置称为「建木」或「天梯」，把母船留在正面的基地称为「广寒宫」，把「昆仑山号」与「玉山」并入「昆仑」神话，甚至把降落时刻的火光编入「天命玄鸟，降而生商」的古老颂歌之中。

这里必须顺带划清一条人口边界：仰韶、龙山与二里头原有的缔造者，以及公元前二二一年一统六国的秦始皇本人，全都是地道的黄河流域本土居民。分六批陆续降落的总人数合计不过十二万，散入数以百万计的中原土著之中，既不足以单独建国，也从未在人口与血缘上取代过本土民众；真正从天上落进中原的，只是早期那一批金属残片、历算度量与营建尺度而已。

\section{零散单证之间的空档}

通盘回检这一章所用到的物证，实际上只有 \texttt{ФВ-1970-18} 登船核对册、\texttt{ФВ-1968-29} 热流遥测备忘、\texttt{ФВ-1963-07} 与 \texttt{ФВ-1974-12} 两批金属残片、\texttt{ФВ-1971-22} 二里头柱洞记录，以及第 117 全宗的几页曲线与清单。把这些分散在不同地层与档案馆里的单点物证串成一套「链路衰减、方案 A 再入、次生黄土沉降、改造旧址并截断天梯」的连续经过，依靠的是十余项互不相干的异常同时压在公元前三世纪前后的黄河流域这一统计聚合。如果逐条拆开来看，鄂尔多斯的 \textsc{ta}-180 异常或许能另有地外撞击解释，二里头的烧结柱洞或许能勉强算作某种特异窑址，朱书「天」「梯」二字也可以当作古人随手涂划的符箓；但要把所有异常全算成彼此无关的巧合，反倒需要替每一处单独编一套额外理由，比接受一次系统性着陆还更牵强。

即便如此，这一章在实证层面仍留着三道填不上的空档：十二个乘员舱各自在豫西与关中的精确落点坐标，至今没有一份原始清单存世；表 \ref{tab:evacuation} 里那八千名不见去向的人口与多出的一千名未入册随员，依然无从找齐；至于约公元前四五五年至前二二一年陆续落地的十二万人，如何在通婚融居与记忆消除场的双重作用下于四代之内消退天外认同、又怎样将血脉延绵至后世，则留给第十二章去收束。\footnote{关于落地后人群如何受记忆消除场作用而在代际间褪去天外认同，参见索科洛夫《记忆消除场通论》。}
